\label{triptych:brief-synthesis:next}
\section*{Source-Grounded Synthesis Across the Propers}

\subsection*{Grace precedes the economy of scarcity}
Isaiah speaks to those who cannot pay; Psalm 145 looks to the open hand of God; Matthew begins not with the crowd's resources but with Christ's compassion. This convergence is strongest where the Lectionary itself correlates Isaiah and the Gospel. Divine generosity does not deny limited loaves, late hours, or human need. It makes scarcity answerable to a giver whose action precedes the disciples' calculation.

\subsection*{Hearing and feeding belong together without collapsing}
Isaiah orders the thirsty to hear so that they may live. The acclamation teaches that bread alone is insufficient. Matthew nevertheless narrates actual healing and eating. The relationship therefore resists both material reduction and spiritual evasion. God's word gives the meaning and horizon of life; compassion supplies food for bodies. Neither can excuse neglect of the other.

\subsection*{Christ gives through servants}
The disciples do not produce abundance, yet the crowd receives through their hands. Chrysostom's direct reception emphasizes this pedagogy: Christ reveals insufficiency, takes the gift, and makes the disciples ministers of his provision. The offerings prayer supplies a liturgical analogue rather than historical proof of selection—God sanctifies gifts and makes the offerers an enduring offering.

\subsection*{Love stronger than affliction}
Romans stands independently of the Isaiah--Matthew correlation. It contributes no theory that famine or peril is unreal; rather, no suffering or cosmic power can defeat God's love in Christ. Joined cautiously to the final prayer, it directs present refreshment toward enduring protection and eternal redemption. Hope is not confidence in uninterrupted supply but confidence that no creature can finally sever communion with God.

\begin{fourstable}{Relation}{Elements}{Classification}{Limit}
Free nourishment & Isaiah; Gospel & Officially correlated & The relation does not turn grace into a prosperity guarantee.\\
Open hand and feeding & Psalm; Gospel & Responsorial and textual observation & Psalmic kingship includes justice for the vulnerable, not consumption alone.\\
Word and bread & Acclamation; Gospel & Acclamatory and source-grounded synthesis & Bodily hunger remains morally serious.\\
Suffering and communion & Romans; final prayer & Semi-continuous plus source-grounded synthesis & No designed cycle-specific relation is claimed.\\
\end{fourstable}
